% authority.tex -- Methodology 3.5: authority and safety as runtime assurance (drafted 12 Oct 2026 from the thesis
% plan of 11 Oct). The safety paragraph moved here from system_overview.tex (its label sec:safety comes with it);
% fig:owners (F5) replaces the old nodes and cascade figures. The rig switch states in brackets follow plan D6 and
% the M1 plan's D6(a)/(b): they are set per lab day and are not results.

\subsection{Authority and safety}
\label{sec:safety}

Every UAV in the air has exactly one controller sending it commands at every instant, and every phase has a way
down that ends with the UAV disarmed on the floor (Figure~\ref{fig:owners}, Appendix~\ref{sec:app_method}). The planner owns the ring; a tracker owns
its UAV's radio link; a UAV off the ring is owned by its own tracker or by the collaborator's controller, and the
multiplexer in front of the newcomer's radio is the one place where ownership changes hands, latched at the
confirmation. The operator holds a disarm button that overrides everything.



\paragraph{The envelope and the fall-backs.} Every tracker checks its own UAV against a flight envelope (UAV and
ring tilt 60$^\circ$, speed 3\,m/s) and checks that its motion-capture pose (0.25\,s) and the planner's reference
(1\,s) are fresh. A violation on a carrier disarms the whole fleet: one UAV losing thrust on a shared rigid payload drags
the others with it, and in a 6\,$\times$\,6\,$\times$\,4\,m cage a falling ring is a short fall onto a padded floor,
where an armed UAV flying on a wrong reference is not; there is no geofence. The planner and the admission chain are
the high-performance controllers of a Simplex-style arrangement (Section~\ref{sec:lit_switching}), and each fall-back
is simple, has its own monitor and ends with every UAV owned and down: the \textbf{take-back} of a newcomer whose catch
is not confirmed; the \textbf{tilt cut} of a newcomer tilting while it is not carrying the ring; \textbf{LAND}, which
lands every UAV whatever it is doing; \textbf{landing the fleet on a newcomer fault}, with the newcomer disarmed alone;
the \textbf{drop rule}, which releases every magnet and lands every UAV when a loss leaves fewer than three carriers or
the ring tilts past 30$^\circ$, on in every mode since 14 October and fired on the rig on 15 October
(Section~\ref{sec:res_drop}); and the \textbf{freed UAV disarmed alone} at its rest.

\paragraph{What is tested where.} ``A tested way down in every phase'' holds in the twin, where every fall-back
above has been flown from the phase it guards, the attach phases on the earlier attach sequence
(Appendix~\ref{sec:app_faults}). On the rig it holds only for the
exits flown there. The drop rule has fired on the rig, after an uncommanded magnet release in a hover, and every
carrier then landed by itself and was disarmed alone; LAND and the take-back have also flown there; the tilt cut and the flight controller's own
disarm-to-land stay off on the rig until each has passed its review, so until then a take-back that sticks ends in
the operator's disarm. When a ring drop comes while the collaborator's controller flies the newcomer, the rig takes
the newcomer back and lands it (twin-tested), and the operator can ask the collaborator's planner to land it straight
down. Redistribution after a release has flown on the rig; the policy for a true fault on a carrier stays the fleet
stop and the drop rule.
